% Supplement subsections typed below as "Section~S<n>.<m>"; build.sh checks each number against the built supplement.
% supplement-ref supp:gaussian=S1.3
% supplement-ref supp:alphabet=S1.5
% supplement-ref supp:angle=S1.6
% supplement-ref supp:margin=S1.7
% supplement-ref supp:borda=S1.8
% supplement-ref supp:persistent=S1.9
\section{Theoretical Analysis}\label{sec:theory}

What do the fixed parts of ArrowFlow guarantee? This section answers four questions.
\begin{enumerate}[leftmargin=2em,topsep=2pt,itemsep=1pt]
\item When does the encoded ranking stay the same?
\item Which rankings can the encoder produce?
\item How far can an input move before a fixed layer changes its output?
\item What does one batch update compute?
\end{enumerate}
Each statement concerns a fixed encoder, layer or batch. None of them is a convergence or generalization theorem for the trained network. Section~S1 states and proves every result that we summarize here in words.

\subsection{Stability of the Ordinal Encoding}\label{sec:theory-stability}

The encoded ranking $\pi_x=\argsort(z)$ changes only when two coordinates of the score vector $z\in\R^e$ (Section~\ref{sec:encoding}) cross. So the smallest gap between two scores, $\delta_{\min}(z)=\min_{i<j}|z_i-z_j|$, controls how stable the ranking is.

\begin{proposition}[Score-gap stability]\label{prop:score-stability}
Let $z\in\R^e$, $e\ge2$, have distinct coordinates, and let $\varepsilon\in\R^e$ satisfy $\|\varepsilon\|_\infty<\delta_{\min}(z)/2$. Then $\argsort(z+\varepsilon)=\argsort(z)$.
\end{proposition}

In words, if every score moves by less than half the smallest gap, the ranking stays the same. The condition is sufficient but not necessary. Its bound cannot be raised, because the two closest scores cross when each moves toward the other by more than half their gap (Figure~\ref{fig:cones}b). Nor does it give a noise threshold, because Gaussian noise exceeds any bound with positive probability. For Gaussian input noise and an encoder of degree one, Section~S1.3 bounds the probability that the ranking changes.

\subsection{The Ordinal Alphabet}\label{sec:theory-alphabet}

Ranking keeps order and discards magnitude, so the encoder can output at most $e!$ rankings. All score vectors without ties that share one ranking $\pi$ form an open convex cone, $\{z: z_{\pi[1]}<z_{\pi[2]}<\dots<z_{\pi[e]}\}$, cut out by the hyperplanes $z_i=z_j$ (Figure~\ref{fig:cones}). A vector with tied scores lies on a wall of a cone and gets the ranking that the ID rule assigns. The encoding is constant inside a cone and jumps at its walls. A fixed projection may land in only some of these cones. So $e!$ bounds the alphabet of rankings that the encoder can emit, but it does not measure the information the data carry (Section~S1.5).

At degree one the alphabet can be counted. Every score is then the dot product of the $d$ standardized features with a fixed vector of the view, a column of $W$ for a random view. So the ranking orders $e$ fixed points in $d$ dimensions by their projection onto one direction. \citet{cover1967orderings} counted how many orderings such projections can produce. A view of degree one reaches at most that many rankings, and a random view can reach all of them. For $e=16$ and $d=4$ there are 437,040, against $16!\approx2.1\times10^{13}$ (Section~S1.5). Polynomial expansion makes the scores linear in more coordinates, so it can enlarge this set.

For the random projection, inputs that point in similar directions get similar rankings. Suppose that the expanded, standardized features of two fixed inputs form the angle $\theta$. Then, in expectation over the random projection matrix, their rankings disagree on $\binom{e}{2}\theta/\pi$ coordinate pairs (Section~S1.6).

\begin{figure}[!htbp]
\centering
\resizebox{0.55\textwidth}{!}{%
\begin{tikzpicture}[scale=1.0]
\node[font=\small, anchor=west] at (-4.6, 3.9) {(a)};
\node[font=\scriptsize, anchor=west] at (-3.9, 3.9) {section $z_1{+}z_2{+}z_3=0$};
% The three walls {z_i = z_j} all contain the diagonal z_1=z_2=z_3, so in a transverse
% section they are three concurrent lines through the origin, 60 degrees apart, and they
% cut the section into exactly six sectors, one per ordering of three coordinates.
% one dash style per wall, so that the walls differ by more than color
\draw[dashed, thick, blue!70] (0, -4.2) -- (0, 4.2) node[above, font=\scriptsize, blue!80!black] {$z_1{=}z_2$};
\draw[dash dot, thick, green!50!black] (-3.637, -2.1) -- (3.637, 2.1) node[right, font=\scriptsize, green!50!black] {$z_2{=}z_3$};
\draw[dash dot dot, thick, red!70] (3.637, -2.1) -- (-3.637, 2.1) node[left, font=\scriptsize, red!80!black] {$z_1{=}z_3$};
\fill[black!45] (0,0) circle (1.6pt);
% Sector labels, counter-clockwise from the 0 degree direction
\node[font=\scriptsize, fill=white, inner sep=2pt, rounded corners=2pt] at (4.5, -0.55) {$z_1{>}z_3{>}z_2$};
\node[font=\scriptsize, fill=white, inner sep=2pt, rounded corners=2pt] at (1.45, 2.512) {$z_1{>}z_2{>}z_3$};
\node[font=\scriptsize, fill=white, inner sep=2pt, rounded corners=2pt] at (-1.45, 2.512) {$z_2{>}z_1{>}z_3$};
\node[font=\scriptsize, fill=white, inner sep=2pt, rounded corners=2pt] at (-2.9, 0) {$z_2{>}z_3{>}z_1$};
\node[font=\scriptsize, fill=white, inner sep=2pt, rounded corners=2pt] at (-1.45, -2.512) {$z_3{>}z_2{>}z_1$};
\node[font=\scriptsize, fill=white, inner sep=2pt, rounded corners=2pt] at (1.45, -2.512) {$z_3{>}z_1{>}z_2$};
% Data point and stability circle, drawn to scale. The centered point (1,-1,0) sits on the
% 0 degree ray at distance sqrt(2) from the origin; its nearest wall is at 1/sqrt(2), so with
% delta_min(z) = 1 the certified radius 1/2 is 0.707 of that distance.
\fill[black] (2.6, 0) circle (2pt);
\draw[thick, dotted, black!40] (2.6, 0) circle (0.919);
\draw[black!50] (2.6, 0) -- (1.681, 0);
\node[font=\scriptsize, black!75, anchor=north] at (2.14, -0.07) {$\delta_{\min}(z)/2$};
\node[font=\scriptsize, anchor=west] at (3.6, 1.1) {$z = (3, 1, 2)$};
\node[font=\scriptsize, anchor=west, black!75] at (3.6, 0.7) {$\argsort(z) = [2,3,1]$};
\end{tikzpicture}%
}\\[8pt]
\begin{tikzpicture}[>=Stealth, lbl/.style={font=\scriptsize}]
\node[font=\small, anchor=west] at (-0.6, 1.6) {(b)};
% one number line of e projected scores; shaded: the certified radius delta_min/2 around each score
\node[lbl, anchor=west] at (0.0, 1.6) {projected scores $z_1, \dots, z_e$ on one line; each shaded band has half-width $\delta_{\min}(z)/2$};
\foreach \x in {0.9, 3.0, 5.2, 6.2, 8.9, 11.9} {
    \fill[black!10] (\x-0.5, 0.48) rectangle (\x+0.5, 0.72);
}
\draw[thick] (0.0, 0.6) -- (13.2, 0.6);
\foreach \x in {0.9, 3.0, 5.2, 6.2, 8.9, 11.9} {
    \fill (\x, 0.6) circle (2.2pt);
}
% a perturbation inside the radius: every score stays in its band, the ranking is unchanged
\node[lbl, blue!70, anchor=west] at (0.0, 1.28) {every score moves less than $\delta_{\min}(z)/2$: the ranking is unchanged};
\foreach \x/\d in {0.9/0.35, 3.0/-0.3, 5.2/0.35, 6.2/-0.35, 8.9/0.3, 11.9/-0.35} {
    \draw[->, thick, blue!70] (\x, 0.98) -- (\x+\d, 0.98);
}
% the smallest gap
\draw[<->, thick, black!50] (5.2, 0.3) -- (6.2, 0.3);
\node[lbl, anchor=north] at (5.7, 0.26) {$\delta_{\min}(z)$};
% a perturbation beyond the radius: two neighbors cross, the ranking changes
\draw[->, thick, red!70] (5.2, -0.36) -- (5.9, -0.36);
\draw[->, thick, red!70] (6.2, -0.52) -- (5.5, -0.52);
\node[lbl, red!70, anchor=west, align=left] at (6.6, -0.44) {two neighbors move more than $\delta_{\min}(z)/2$ and cross:\\ the ranking changes};
\end{tikzpicture}
\caption{\textbf{Permutation cones and the stability radius.} (a)~Up to the hyperplanes $\{z_i=z_j\}$, whose points the tie rule assigns, the argsort encoding partitions the projected space $\R^e$ into $e!$ open convex cones (Weyl chambers). They are drawn here for $e=3$ in the plane $z_1+z_2+z_3=0$, which every cone meets and which the three hyperplanes cut into six sectors. All points in a cone map to the same permutation. The point $z=(3,1,2)$ is drawn at its position in the plane, which differs from $z$ by a multiple of the all-ones vector and has the same ranking and the same coordinate gaps. The dotted circle has radius $\delta_{\min}(z)/2$, half the smallest coordinate gap of Proposition~\ref{prop:score-stability}. It lies inside the region $\|\varepsilon\|_\infty<\delta_{\min}(z)/2$, so a perturbation inside it cannot cross a boundary. The nearest wall is farther away, at distance $\delta_{\min}(z)/\sqrt2$. (b)~The same condition on one line of projected scores. Scores that each move by less than $\delta_{\min}(z)/2$ keep their ranking, and two neighbors that move further can cross.}
\label{fig:cones}
\end{figure}


\subsection{Margins and Ties of a Fixed Ranking Layer}\label{sec:theory-margin}

A fixed layer is stable in the same way. The footrule is a metric, so no filter's distance changes by more than the input moves. When all distances differ, the smallest gap between consecutive distances therefore guarantees the same output for every input within half that gap (Section~S1.7). The guarantee fails where distances tie. Footrule distances are even integers up to $\lfloor V^2/2\rfloor$ \citep{diaconis1977footrule}, so at most $\lfloor V^2/4\rfloor+1$ values occur. A layer with more filters than that has tied distances at every input. For example, a layer of $128$ filters over $V=16$ items can see at most $65$ distinct distances. Ties are common well before that bound. The distances from an input to random filters crowd around their mean $(V^2-1)/3$ with a standard deviation near $0.21\,V^{3/2}$. At the hidden-layer sizes used here, random filters therefore leave almost no input with all distances distinct, and the guarantee then covers at most the nearest filter (Section~S1.7). Where distances tie, the tie rule of Section~\ref{sec:forward} decides the output order.

\subsection{What One Batch Update Computes}\label{sec:theory-aggregation}

One batch update is a weighted Borda count of the accepted votes and the prior. In general it is not the footrule median. Because the accumulator is reset after every batch, it keeps no running estimate of where each item belongs on average. Each update depends only on the filter's current order and the current batch. With little vote weight, the current order can even hold against the average vote (Section~S1.9).

\begin{proposition}[One batch update is a weighted Borda count]\label{prop:batch-borda}
Fix a filter $r_j$ and the targets $\tilde\tau_1,\dots,\tilde\tau_n$ of one batch with vote weights $|a_1|,\dots,|a_n|>0$. After accumulation, the mean voted position $s_p$ of~\eqref{eq:rowmean} for the row $p=\operatorname{pos}_{r_j}(v)$ of item $v$ equals
\begin{equation}\label{eq:rowmean-borda}
s_p=\frac{\operatorname{pos}_{r_j}(v)+\sum_{t=1}^{n}|a_t|\operatorname{pos}_{\tilde\tau_t}(v)}{1+\sum_{t=1}^{n}|a_t|}=\frac{P_j(v)}{1+\sum_{t=1}^{n}|a_t|},
\end{equation}
so in exact arithmetic the reordered filter is the ranking by $P_j$ of~\eqref{eq:borda-update} with the fixed tie rule. It differs in general from the footrule median of the same profile, even when the mean positions are distinct.
\end{proposition}

The weighted footrule median minimizes the summed absolute position differences to the voters, and minimum-cost matching computes it exactly \citep{dwork2001rank}. The Borda ranking instead minimizes the summed squared differences \citep{zwicker2008average}. For equal-weight votes, it is also the maximum-likelihood center estimate in a ranking model based on this squared distance \citep{fligner1986distance} (Section~S1.8). The profiles of Section~S1.8 show that the two rules do disagree, on a finite profile and in the limit. Picture an idealized estimator that kept every vote drawn from one fixed distribution. It would converge to the ranking by expected position when the expected positions differ. The reset update need not converge to it, and in the network the votes also change with the current predictions (Section~S1.9).
